% !PRINTED 288
\weberpage{288}
%<<CONT
\weq{16}{\begin{pmatrix} \alpha, & \beta \\ \gamma, & \delta \end{pmatrix}, \quad \alpha\delta - \beta\gamma \equiv 1 \ (\mathrm{mod}\ p).}
In one of these substitutions the numbers $\alpha$, $\beta$, $\gamma$, $\delta$ are determined, with respect to the modulus $p$, up to the common sign. The degree of the group $\mathfrak{L}_0$ is
\weq{17}{\frac{1}{2}p(p^2-1).}

One arrives at the form (16) directly, however, if one takes as the group of the division equation not the group $\mathfrak{A}$ but the group $\mathfrak{B}$ of § 63, that is, if one adjoins \gesperrt{the $p$th roots of unity to the domain of rationality}; from which the theorem flows:

\gesperrt{The group $\mathfrak{L}_0$ is the group of the transformation equation when the $p$th roots of unity are adjoined to the domain of rationality.}

For the reduction of the group $\mathfrak{L}$ to the group $\mathfrak{L}_0$, however, the adjunction of a two-valued function already suffices, and with the adjunction of the $p$th roots of unity too much has therefore been done. To recognise which irrationality it is \gesperrt{necessary} to adjoin, the theorems of §§ 63, 65 serve.

In § 63, 3. we have seen that the $p$th root of unity $\varrho$ is rationally representable by the roots of the division equation, and that by a substitution of the group $\mathfrak{A}$ whose determinant is congruent to $m$, $\varrho$ passes over into $\varrho^m$; further, we proved in § 65 that by the adjunction of all the roots of a transformation equation the group of the division equation is reduced to the group $\mathfrak{A}_0$, which consists of all substitutions of the form
\weq{18}{\begin{pmatrix} a, & 0 \\ 0, & a \end{pmatrix}}
where $a$ is an arbitrary number not divisible by $p$. The determinants of the substitutions of $\mathfrak{A}_0$ are therefore squares, and are consequently each congruent modulo $p$ to one of the $\frac{1}{2}(p-1)$ \gesperrt{quadratic residues} of $p$.

The sum
\weq{19}{A = \overset{a}{\Sigma}\varrho^a,}
in which all the \gesperrt{quadratic residues} of $p$ are to be put for $a$, therefore remains unchanged by the substitutions
%CONT>>
